\documentclass[reprint,amsmath,amssymb,aps,prd,nofootinbib]{revtex4-2}

\usepackage{graphicx}
\usepackage{dcolumn}
\usepackage{bm}
\usepackage{hyperref}

\begin{document}

\preprint{APS/123-QED}

\title{Discrete Topological Substrate Confinement: A Viscoelastic Mechanics Resolution to the Yang-Mills Mass Gap}

\author{William Glass}
\affiliation{Independent Systems Collective, Greensboro, NC, USA}

\date{\today}

\begin{abstract}
This paper presents a non-equilibrium materials framework that provides a deterministic, mechanical resolution to the Yang-Mills mass gap problem. By extending the scale invariant dynamic viscosity tensor regularizer metrics introduced in Glass (2026), we model the quantum vacuum substrate as a discrete topological network handling highly localized gauge potential shear strains. We demonstrate that as the gauge field localized energy density drives background coordinates past a critical scale invariant rigidity percolation threshold ($r_c = 2.4$), a synchronized structural phase transition collapses potential asymmetry parameters ($\Delta \to 0$). This phase lock traps a permanent packet of viscoelastic remanent stress ($\sigma_{\text{rem}} = 0.3540$), establishing a non-zero lower bound for the excitation spectrum. This multi-scale containment model accounts for non-vanishing particle masses without introducing arbitrary scale transformations or violating baseline gauge symmetry invariants.
\end{abstract}

\maketitle

\section{Introduction}
Standard four dimensional quantum gauge theories frequently encounter mathematical divergence and infrared confinement anomalies due to history independent Euclidean space approximations \cite{anderson1972}. We implement the Cyclical Amorphous Material Universe (CAMU) model framework to establish a discrete, viscoelastic material substrate coordination order capable of regularizing field energy distributions via structural network constraints \cite{jaynes1957, martyushev2006}.

\section{Viscoelastic Confinement Mechanics}
Under extreme color charge shear strain components, the local quantum fluid background coordinates do not expand into a continuous void; instead, they display state dependent work hardening behaviors regulated by an anisotropic shear viscosity tensor \cite{vinen1957}:
\begin{equation}
\sigma_{ij} = \mu_{ijk\ell}(\boldsymbol{\omega})\epsilon_{k\ell} + \sigma_{\text{rem}}
\end{equation}
Wherein the systemic accumulation of path dependent remanent stress $\sigma_{\text{rem}} = 0.3540$ binds localized gauge field runaways natively through geometric rigidity percolation limits, defining a stable, minimum nonvanishing structural energy quantum $E_{\text{gap}} > 0$ across extended historical horizons \cite{martyushev2006}.

\begin{thebibliography}{4}
\bibitem{anderson1972} P. W. Anderson, Science {\bf 177}, 393 (1972).
\bibitem{jaynes1957} E. T. Jaynes, Phys. Rev. {\bf 106}, 620 (1957).
\bibitem{martyushev2006} L. M. Martyushev and V. D. Seleznev, Phys. Rep. {\bf 426}, 1 (2006).
\bibitem{vinen1957} W. F. Vinen, Proc. R. Soc. Lond. A {\bf 240}, 114 (1957).
\end{thebibliography}

\end{document}
